\chapter{A-TP via Nyman--Beurling Closure and the Baez-Duarte Refinement}
\label{v4-arith:w7-b:chapter}

\emph{The Nyman--Beurling comparison. Chapter~\ref{v4-arith:w6-a:ATP-direct-obstruction} reduced the archimedean
Tate positivity conjecture A-TP to a bounded-interval question on
\([-t_{\ast},t_{\ast}]\) with \(t_{\ast}\approx 0.8507\); the irreducible
residual is the elementary ceiling. The same obstruction has a
classical \(L^{2}(0,1)\) expression through Nyman 1950 and
Beurling 1955, sharpened by Baez-Duarte 2003. The Nyman--Beurling
criterion reformulates RH as a closure statement in \(L^{2}(0,1)\)
generated by the functions \(\rho_{\alpha}(x)=\{\alpha/x\}-\alpha\{1/x\}\)
for \(\alpha\in(0,1]\); Baez-Duarte shows RH is equivalent to the
asymptotic vanishing of the M\"{o}bius-coefficient distance
\(d_{N}=\inf_{c_{k}}\|\chi_{[0,1]}-\sum_{k=1}^{N}c_{k}\rho_{1/k}\|_{L^{2}(0,1)}\).
The results are fourfold. First, an explicit
Mellin normal form: the \(L^{2}(0,1)\) norm of
\(\chi_{[0,1]}-\sum c_{k}\rho_{1/k}\) is an integral on the critical
line against a kernel built from \(\zeta(s)\).  A-TP and the
Baez-Duarte vanishing have the same RH truth value; the Mellin identity
does not identify the finite-dimensional distance with the Weil
quadratic form.
Second, on the bounded-interval reduction HFD sub-class of
Chapter~\ref{v4-arith:w6-a:ATP-direct-obstruction}, the expected
Nyman--Beurling estimate becomes a weighted approximation problem for
\(\zeta(s)^{-1}\) on the critical line; neither the Mellin identity nor
PNT proves the \(1/\log N\) rate. Third, a M\"{o}bius partial-sum
control: the conditional Mertens bound
\(M(N):=\sum_{k\leq N}\mu(k)=O(\sqrt{N}\log^{2}N)\) under RH
sharpens the unconditional PNT bound \(M(N)=o(N)\); the
corresponding \(M_{1}(N):=\sum_{k\leq N}\mu(k)/k=o(1)\)
(unconditional, PNT) is only a scalar shadow of the critical-line
\(L^{2}\)-approximation. Fourth, the complementary class carries a
compact-core Baez-Duarte shadow of the same bounded-interval residual;
this compact shadow is not a substitute for the global
Nyman--Beurling criterion.}

\section{Primary Sources}
\label{v4-arith:w7-b:sources}

The four source lists below are disjoint from the Weil 1952 /
Rankin--Selberg / Petersson lists driving
Chapters~\ref{v4-arith:w5-a:chapter} and
\ref{v4-arith:w6-a:ATP-direct-obstruction}, and disjoint from the de
Branges / Hermite--Biehler list of
Chapter~\ref{v4-arith:w5-j:VBstar-spectral-positivity}.

\begin{description}
\item[Nyman 1950.] Nyman, B. \emph{On the One-Dimensional Translation
Group and Semi-Group in Certain Function Spaces}. PhD thesis, Uppsala
1950. Introduces the closed subspace \(H^{2}(0,1)\) of \(L^{2}(0,1)\)
generated by the translated fractional-part functions and proves the
\(\chi_{[0,1]}\in\overline{H^{2}(0,1)}\) criterion for RH.
\item[Beurling 1955.] Beurling, A. \emph{A closure problem related
to the Riemann zeta-function}. \emph{Proc.~Nat.~Acad.~Sci.~USA} 41
(1955) 312--314. Reformulates Nyman's criterion in the modern form
with the generators \(\rho_{\alpha}(x)=\{\alpha/x\}-\alpha\{1/x\}\) and
establishes the equivalence \(\mathrm{RH}\Leftrightarrow \chi_{[0,1]}\in
H^{2}(0,1)\).
\item[Baez-Duarte 2003.] Baez-Duarte, L. \emph{A strengthening of
the Nyman--Beurling criterion for the Riemann hypothesis}.
\emph{Atti~Accad.~Naz.~Lincei Cl.~Sci.~Fis.~Mat.~Natur.~Rend.~Lincei
(9) Mat.~Appl.} 14 (2003) 5--11, arXiv:math/0202141. Restriction to
\(\alpha=1/k\) for \(k\in\mathbb{N}\); M\"{o}bius coefficients
\(c_{k}=\mu(k)/k\); distance \(d_{N}\to 0\) equivalent to RH.
\item[Baez-Duarte--Balazard--Landreau--Saias 2005.] Baez-Duarte, L.;
Balazard, M.; Landreau, B.; Saias, E. \emph{Notes sur la fonction
\(\zeta\) de Riemann}, Part IV. \emph{Adv.~Math.} 188 (2005) 69--86.
They give the \(1/\log N\) lower obstruction:
\(\liminf d_N^2\log N\geq \sum_\rho m_\rho^2/|\rho|^2\).
\item[Burnol 2001.] Burnol, J.-F. \emph{A lower bound in an approximation
problem involving the zeros of the Riemann zeta function},
arXiv:math/0103058. The multiplicity \(m_\rho^2\) enters the lower
bound through the local principal part at \(\rho\).
\end{description}

Independence of inputs: Nyman 1950, Beurling 1955, Baez-Duarte 2003,
and Baez-Duarte--Balazard--Landreau--Saias 2005 form a
Uppsala--Caracas--Bordeaux line of \(L^{2}(0,1)\) closure analysis, none
entering the Weil 1952 explicit-formula derivation, the Rankin--Selberg
identity, the Jacquet--Langlands adelic formalism, or the de Branges
Hermite--Biehler apparatus. The Vol IV disjointness check against
Chapters~\ref{v4-arith:w5-a:chapter},
\ref{v4-arith:w6-a:ATP-direct-obstruction},
\ref{v4-arith:w5-j:VBstar-spectral-positivity},
\ref{v4-arith:w5-i:chapter} passes on a per-source basis.

\section{The Nyman--Beurling Space \(H^{2}(0,1)\)}
\label{v4-arith:w7-b:NB-space}

\subsection{Generators and Mellin Transform}
\label{v4-arith:w7-b:NB-space:generators}

\begin{definition}[Nyman--Beurling generators]
\label{v4-arith:w7-b:def:rho-alpha}
\ClaimStatusDefinitional. For \(\alpha\in(0,1]\) define
\(\rho_{\alpha}\colon(0,1)\to\mathbb{R}\) by
\begin{equation}
\label{v4-arith:w7-b:eq:rho-alpha}
\rho_{\alpha}(x)\;:=\;\left\{\frac{\alpha}{x}\right\}-\alpha\!\left\{\frac{1}{x}\right\},\qquad x\in(0,1),
\end{equation}
where \(\{y\}=y-\lfloor y\rfloor\) is the fractional part. Each
\(\rho_{\alpha}\in L^{2}(0,1)\) (Nyman 1950, Beurling 1955). The
Nyman--Beurling closed subspace is
\begin{equation}
\label{v4-arith:w7-b:eq:H2-space}
H^{2}(0,1)\;:=\;\overline{\mathrm{span}\{\rho_{\alpha}:\alpha\in(0,1]\}}\;\subset\;L^{2}(0,1),
\end{equation}
the closed linear span in \(L^{2}\)-norm.
\end{definition}

\begin{lemma}[Mellin transform of \(\rho_{\alpha}\)]
\label{v4-arith:w7-b:lem:mellin-rho}
\ClaimStatusProvedElsewhere. For \(\mathrm{Re}\,s>0\), \(s\neq 1\),
\begin{equation}
\label{v4-arith:w7-b:eq:mellin-rho}
\int_{0}^{1}\rho_{\alpha}(x)\,x^{s-1}\,dx
\;=\;\frac{\alpha-\alpha^{s}}{s}\,\zeta(s),
\end{equation}
by analytic continuation from \(\mathrm{Re}\,s>1\). The zero of
\(\alpha-\alpha^s\) at \(s=1\) cancels the simple pole of \(\zeta\).
\end{lemma}

\begin{proof}
Beurling 1955, Equation (3); Baez-Duarte 2003 Lemma~2.1. The
calculation is elementary: \(\int_{0}^{1}\{\alpha/x\}x^{s-1}dx=
\alpha^{s}\int_{\alpha}^{\infty}\{y\}y^{-s-1}dy\) via the substitution
\(y=\alpha/x\), and the fractional-part Dirichlet series
\(\int_{1}^{\infty}\{y\}y^{-s-1}dy=1/(s-1)-\zeta(s)/s\) (classical;
Titchmarsh \S2.1). Since \(0<\alpha\le1\),
\[
\int_{\alpha}^{1} y\,y^{-s-1}\,dy=\frac{1-\alpha^{1-s}}{1-s}.
\]
Hence
\[
\int_{0}^{1}\{\alpha/x\}x^{s-1}dx
=\alpha^{s}\left(\frac{1-\alpha^{1-s}}{1-s}
+\frac{1}{s-1}-\frac{\zeta(s)}{s}\right)
=\frac{\alpha}{s-1}-\frac{\alpha^s\zeta(s)}{s}.
\]
Subtracting \(\alpha\int_0^1\{1/x\}x^{s-1}dx
=\alpha(1/(s-1)-\zeta(s)/s)\) gives
\begin{align*}
\int_{0}^{1}\rho_{\alpha}(x)x^{s-1}dx
&=\left(\frac{\alpha}{s-1}-\frac{\alpha^s\zeta(s)}{s}\right)
-\alpha\left(\frac{1}{s-1}-\frac{\zeta(s)}{s}\right)\\
&=\frac{\alpha-\alpha^s}{s}\zeta(s).
\end{align*}
\end{proof}

\begin{definition}[the characteristic function]
\label{v4-arith:w7-b:def:chi}
\ClaimStatusDefinitional. Set \(\chi(x):=\chi_{[0,1]}(x)\) restricted to
\(L^{2}(0,1)\), i.e.\ the constant function \(1\) on \((0,1)\). Its Mellin
transform is
\begin{equation}
\label{v4-arith:w7-b:eq:mellin-chi}
\int_{0}^{1}\chi(x)x^{s-1}dx=\frac{1}{s},\qquad\mathrm{Re}\,s>0.
\end{equation}
\end{definition}

\subsection{The Nyman--Beurling Theorem (Recall)}
\label{v4-arith:w7-b:NB-space:theorem}

\begin{theorem}[Nyman--Beurling 1950/1955]
\label{v4-arith:w7-b:thm:nyman-beurling}
\ClaimStatusProvedElsewhere. The Riemann hypothesis is equivalent to
\begin{equation}
\label{v4-arith:w7-b:eq:NB-criterion}
\chi\;\in\;H^{2}(0,1).
\end{equation}
That is: \(\chi\) lies in the closed subspace generated by the
\(\{\rho_{\alpha}\}\) if and only if RH holds.
\end{theorem}

\begin{proof}
Nyman 1950 Ch.~III; Beurling 1955. The direction
RH \(\Rightarrow\chi\in H^{2}\) uses explicit approximants by rational
sums; the converse uses that \(\zeta^{-1}\) extending to the strip
\(\mathrm{Re}\,s\geq 1/2+\varepsilon\) controls the closure.
\end{proof}

\begin{theorem}[Baez-Duarte refinement]
\label{v4-arith:w7-b:thm:baez-duarte}
\ClaimStatusProvedElsewhere. RH is equivalent to
\(\lim_{N\to\infty}d_{N}=0\), where
\begin{equation}
\label{v4-arith:w7-b:eq:dN-definition}
d_{N}\;:=\;\inf_{c_{1},\ldots,c_{N}\in\mathbb{C}}\;\bigl\|\chi-\sum_{k=1}^{N}c_{k}\rho_{1/k}\bigr\|_{L^{2}(0,1)}.
\end{equation}
Equivalently: RH holds iff the distance from \(\chi\) to the finite-rank
Nyman--Beurling space
\(H^{2}_{N}:=\mathrm{span}\{\rho_{1/k}:k=1,\ldots,N\}\) tends to zero
as \(N\to\infty\).
\end{theorem}

\begin{proof}
Baez-Duarte 2003 Theorem~1.1. The restriction to \(\alpha=1/k\) and
the passage from the closure to the sequential vanishing is the
content of the refinement.
\end{proof}

\begin{theorem}[Baez-Duarte--Balazard--Landreau--Saias lower bound]
\label{v4-arith:w7-b:thm:bdbls-rate}
\ClaimStatusProvedElsewhere. Let \(m_{\rho}\) be the multiplicity of the
non-trivial zero \(\rho\). Then
\begin{equation}
\label{v4-arith:w7-b:eq:bdbls-rate}
\liminf_{N\to\infty} d_N^2\log N
\;\geq\;
\sum_{\rho}\frac{m_{\rho}^{2}}{|\rho|^{2}}.
\end{equation}
If RH holds and all zeros are simple, the right hand side is
\[
\sum_{\rho}\frac{1}{|\rho|^{2}}
=2+\gamma_{E}-\log(4\pi)=0.04619\ldots .
\]
Thus \(d_N\) cannot be \(o((\log N)^{-1/2})\).  A matching asymptotic
with this leading constant is a stronger theorem and is not used here.
\end{theorem}

\begin{proof}
Baez-Duarte--Balazard--Landreau--Saias 2005 and Burnol's multiplicity
refinement give the lower bound. The evaluation of the simple-zero
constant follows by differentiating the Hadamard product for
\(\xi_{\mathbb R}\) at \(s=0\).
\end{proof}

\section{Mellin Normal Form and RH-Level Comparison}
\label{v4-arith:w7-b:mellin-translation}

The Mellin--Plancherel identity converts the Nyman--Beurling
\(L^{2}(0,1)\) norm into a critical-line integral.  The Weil explicit
formula and the Nyman--Beurling criterion both have RH as their
zero-placement content; this is equality of logical strength, not an
identity of quadratic forms.

\subsection{The Mellin Norm Identity}
\label{v4-arith:w7-b:mellin-translation:norm}

\begin{lemma}[Mellin--Plancherel for \(L^{2}(0,1)\) in the half-line model]
\label{v4-arith:w7-b:lem:mellin-plancherel}
\ClaimStatusProvedHere. For \(\phi\in L^{2}(0,1)\) with Mellin
transform \(\widehat{\phi}_{M}(s):=\int_{0}^{1}\phi(x)x^{s-1}dx\) regular
on \(\mathrm{Re}\,s=1/2\),
\begin{equation}
\label{v4-arith:w7-b:eq:mellin-plancherel}
\|\phi\|_{L^{2}(0,1)}^{2}\;=\;\frac{1}{2\pi}\int_{-\infty}^{\infty}|\widehat{\phi}_{M}(1/2+it)|^{2}\,dt.
\end{equation}
\end{lemma}

\begin{proof}
The standard Mellin--Plancherel theorem (Titchmarsh \emph{The Theory
of Functions} 2nd~ed. \S7.1) gives, for
\(\phi\in L^{2}((0,1),dx)\) extended by zero to
\(L^{2}((0,\infty),dx)\) and the change of variable
\(x=e^{-u}\), an isometry on the line
\(\mathrm{Re}\,s=1/2\). The precise identity
\eqref{v4-arith:w7-b:eq:mellin-plancherel} uses the weighted
\(L^{2}(0,1)\) norm without the \(dx/x\) measure; a substitution
\(x=e^{-u}\) gives
\(\|\phi\|_{L^{2}(0,1)}^{2}=\int_{0}^{\infty}|\phi(e^{-u})|^{2}e^{-u}du\).
For \(\phi=\chi\) constant one,
\(\|\chi\|^{2}=1\) and the Mellin transform \(1/s\) pairs to
\((2\pi)^{-1}\int|1/(1/2+it)|^{2}dt=(2\pi)^{-1}\int dt/(1/4+t^{2})=1\),
matching. For general \(\phi\in L^{2}(0,1)\) the same substitution
reduces to the standard Paley--Wiener/Plancherel identity on the
half-plane. Detailed computation (Baez-Duarte 2003 Lemma~3.1).
\end{proof}

\begin{proposition}[Nyman--Beurling distance as Mellin integral]
\label{v4-arith:w7-b:prop:dN-mellin}
\ClaimStatusProvedHere. For \(c_{1},\ldots,c_{N}\in\mathbb{C}\),
\begin{equation}
\label{v4-arith:w7-b:eq:dN-mellin}
\|\chi-\sum_{k=1}^{N}c_{k}\rho_{1/k}\|_{L^{2}(0,1)}^{2}
=\frac{1}{2\pi}\int_{-\infty}^{\infty}\!\!\left|\frac{1}{s}-\frac{\zeta(s)}{s}\sum_{k=1}^{N}c_{k}\bigl(k^{-1}-k^{-s}\bigr)\right|^{2}_{s=1/2+it}\!dt.
\end{equation}
Equivalently, abbreviating \(A_{N}(s):=\sum_{k=1}^{N}c_{k}\,k^{-s}\),
\begin{equation}
\label{v4-arith:w7-b:eq:dN-mellin-compact}
\|\chi-\sum c_{k}\rho_{1/k}\|_{L^{2}(0,1)}^{2}
=\frac{1}{2\pi}\int_{-\infty}^{\infty}\!\!\left|\frac{1-\zeta(s)\bigl(A_{N}(1)-A_{N}(s)\bigr)}{s}\right|^{2}_{s=1/2+it}\!dt,
\end{equation}
where \(A_{N}(1)=\sum_{k=1}^{N}c_{k}/k\).
\end{proposition}

\begin{proof}
By \Cref{v4-arith:w7-b:lem:mellin-rho},
\(\widehat{\rho_{1/k}}_{M}(s)=\zeta(s)(k^{-1}-k^{-s})/s\).
Linearity: the Mellin transform of \(\chi-\sum c_{k}\rho_{1/k}\) is
\(1/s-\zeta(s)s^{-1}\sum c_{k}(k^{-1}-k^{-s})\). Apply
\Cref{v4-arith:w7-b:lem:mellin-plancherel} on
\(\mathrm{Re}\,s=1/2\).
\end{proof}

\begin{remark}[M\"{o}bius choice and the missing \((s-1)^{-1}\) pole]
\label{v4-arith:w7-b:rem:mobius-cancels}
Setting \(c_{k}=\mu(k)/k\) gives
\(A_{N}(s)=\sum_{k=1}^{N}\mu(k)k^{-s-1}\), so
\(A_{N}(1)=\sum_{k=1}^{N}\mu(k)/k^{2}\to 6/\pi^{2}\) as \(N\to\infty\)
(absolutely convergent, unconditional). There is no
\((s-1)^{-1}\)-term in \eqref{v4-arith:w7-b:eq:dN-mellin-compact}. The
main term is the critical-line factor
\(\bigl(1-\zeta(s)(A_N(1)-A_N(s))\bigr)/s\), and it is not controlled by
the scalar \(M_1(N)=\sum_{k\le N}\mu(k)/k\) alone.
\end{remark}

\subsection{The Weil-Type Insertion}
\label{v4-arith:w7-b:mellin-translation:weil-insertion}

\begin{theorem}[A-TP and Nyman--Beurling have the same RH content]
\label{v4-arith:w7-b:thm:translation-ATP-NB}
\ClaimStatusProvedHere. The archimedean Tate positivity A-TP of
\Cref{v4-arith:w5-a:conj:archimedean-tate-positivity} is equivalent
to the Nyman--Beurling closure \(\chi\in H^{2}(0,1)\); equivalently to
the Baez-Duarte vanishing \(\lim_{N\to\infty}d_{N}=0\).
\end{theorem}

\begin{proof}
Both conditions are equivalent to RH. A-TP
\(\Leftrightarrow\) RH by
\Cref{v4-arith:w5-a:thm:ATP-equals-weil-positivity}. RH
\(\Leftrightarrow\chi\in H^{2}(0,1)\) by
\Cref{v4-arith:w7-b:thm:nyman-beurling}. RH
\(\Leftrightarrow d_{N}\to 0\) by
\Cref{v4-arith:w7-b:thm:baez-duarte}. Transitivity closes the
reformulation. The two equivalent RH formulations are related by an
\emph{explicit} Mellin-Plancherel integral identity.
\end{proof}

\begin{remark}[the Mellin identity is not a zero-square identity]
\label{v4-arith:w7-b:rem:cross-reduction}
Beyond the tautological RH-equivalence, the Nyman--Beurling side gives
the exact Mellin-Plancherel identity
\begin{equation}
\label{v4-arith:w7-b:eq:dN-Weil-link}
\|\chi-\sum_{k=1}^{N}c_{k}\rho_{1/k}\|_{L^{2}(0,1)}^{2}
\;=\;\frac{1}{2\pi}\int_{-\infty}^{\infty}\left|
\frac{1-\zeta(s)\bigl(A_N(1)-A_N(s)\bigr)}{s}
\right|_{s=1/2+it}^{2}\,dt,
\end{equation}
with \(A_N(s)=\sum c_k k^{-s}\). No formula with
\(|\xi_{\mathbb R}(\rho)|^{-2}\) at non-trivial zeros can be correct,
since \(\xi_{\mathbb R}(\rho)=0\). Passing from
\eqref{v4-arith:w7-b:eq:dN-Weil-link} to a zero expansion is exactly the
Baez--Duarte theorem; it is not an elementary residue computation inside
the Mellin identity.
\end{remark}

\section{Restricted-Class Estimate on HFD}
\label{v4-arith:w7-b:restricted-class}

The bounded-interval reduction HFD sub-class (\Cref{v4-arith:w6-a:def:HFD-class})
enters the Nyman--Beurling framework through the spectral weight on the
critical line: the HFD condition concentrates \(|\widehat{f}|^{2}\)
on \(|t|>t_{\ast}\). This concentration gives archimedean
\(W_{\infty}\)-positivity in
\Cref{v4-arith:w6-a:thm:HFD-ATP}; it does not, by itself, give a
weighted Baez-Duarte estimate.

\subsection{HFD-Adapted Nyman--Beurling Distance}
\label{v4-arith:w7-b:restricted-class:HFD}

\begin{definition}[HFD-restricted Nyman--Beurling distance]
\label{v4-arith:w7-b:def:dN-HFD}
\ClaimStatusDefinitional. For \(f\in\mathrm{PW}^{\mathrm{hol},\mathrm{HFD}}_{F}\)
together with a specified unitary Mellin transform
\(\mathcal{M}_{\mathrm{add}\to\mathrm{mult}}\) from the additive
Paley--Wiener model to the critical-line Mellin model, set
\(\widehat f_M=\mathcal{M}_{\mathrm{add}\to\mathrm{mult}}(\widehat f)\)
and define
\begin{equation}
\label{v4-arith:w7-b:eq:dN-HFD}
d_{N}^{\mathrm{HFD}}(f)\;:=\;\inf_{c_{k}}\!\left\|\widehat{f}_{M}(1/2+it)\cdot\!\left[\frac{1-\zeta(s)\sum_{k=1}^{N}c_{k}(k^{-1}-k^{-s})}{s}\right]_{s=1/2+it}\right\|_{L^{2}(\mathbb{R})}.
\end{equation}
This is the Mellin-weighted Nyman--Beurling distance tested against
the spectral profile of \(f\). Without the displayed unitary transform,
\(\widehat f_M\) is not determined by the additive test function.
\end{definition}

\begin{theorem}[weighted HFD Nyman--Beurling bound is an open estimate]
\label{v4-arith:w7-b:thm:HFD-NB-bound}
\ClaimStatusNeedsVerification. The estimate
\begin{equation}
\label{v4-arith:w7-b:eq:HFD-NB-bound}
d_{N}^{\mathrm{HFD}}(f)^{2}
\;\leq\;
C_{\mathrm{HFD}}(f)\cdot\frac{1}{\log N}
\;+\;\varepsilon_{\mathrm{pol}}(N)\cdot\|f\|_{L^{2}}^{2}
\end{equation}
does not follow from the preceding identities. It would require a
uniform weighted Baez--Duarte estimate for the critical-line integral
\eqref{v4-arith:w7-b:eq:dN-Weil-link}. The Mellin identity and the
Prime Number Theorem do not by themselves give this bound.
\end{theorem}

\begin{proof}
The exact identity is \eqref{v4-arith:w7-b:eq:dN-Weil-link}. The
previous zero-expansion contained
\(|\xi_{\mathbb R}(\rho)|^{-2}\), which is undefined at non-trivial
zeros of \(\xi_{\mathbb R}\). What remains is a weighted approximation
problem for \(1/\zeta\) on the critical line. Scalar control of
\(M_1(N)=\sum_{k\leq N}\mu(k)/k\) does not control that weighted
\(L^2\)-norm.
\end{proof}

\begin{corollary}[the HFD Nyman--Beurling estimate is an additional input]
\label{v4-arith:w7-b:cor:HFD-ATP-NB}
\ClaimStatusNeedsVerification. The limit
\[
\lim_{N\to\infty}d_{N}^{\mathrm{HFD}}(f)=0
\]
does not follow from the preceding identities. Thus the identities
do not prove A-TP on HFD.
\end{corollary}

\begin{proof}
It would follow from \Cref{v4-arith:w7-b:thm:HFD-NB-bound} only after
that estimate is known. The exact Mellin identity gives an open
weighted approximation problem.
\end{proof}

\begin{remark}[comparison with the bounded-interval reduction]
\label{v4-arith:w7-b:rem:cross-walk-wave6}
The bounded-interval reduction gives \(W_{\infty}\ge0\) on HFD. The
Nyman--Beurling form adds the critical-line weighted approximation
problem
\eqref{v4-arith:w7-b:eq:HFD-NB-bound}; it does not remove the
residual \(R[f]\), and it does not give full A-TP on HFD without that
extra weighted estimate.
\end{remark}

\section{M\"obius Partial-Sum Control}
\label{v4-arith:w7-b:sec:mobius-control}

The M\"obius partial sums \(M(N)\) and \(M_{1}(N)\) remain useful
background. They do not, by themselves, prove the weighted
HFD estimate of \Cref{v4-arith:w7-b:thm:HFD-NB-bound}.

\begin{proposition}[Mertens-type bounds, unconditional]
\label{v4-arith:w7-b:prop:mertens-uncond}
\ClaimStatusProvedElsewhere. Set \(M(N):=\sum_{k\leq N}\mu(k)\) and
\(M_{1}(N):=\sum_{k\leq N}\mu(k)/k\). Unconditionally:
\begin{align}
\label{v4-arith:w7-b:eq:M-classical}
M(N)&\;=\;o(N),\qquad N\to\infty\quad(\text{PNT, Hadamard-de la Vall\'ee-Poussin 1896}),\\
\label{v4-arith:w7-b:eq:M1-classical}
M_{1}(N)&\;=\;o(1),\qquad N\to\infty\quad(\text{Landau 1911, equivalent to PNT}).
\end{align}
Quantitative form (Titchmarsh \S3.14):
\(M(N)=O(N\exp(-c\sqrt{\log N}))\) and
\(M_{1}(N)=O(\exp(-c\sqrt{\log N}))\) for some \(c>0\).
\end{proposition}

\begin{proof}
Hadamard 1896 (\emph{Bull.~Soc.~Math.~France} 24) and
de la Vall\'ee-Poussin 1896 (\emph{Ann.~Soc.~Sci.~Bruxelles} 20)
established the PNT via zero-free region of \(\zeta\) just inside
\(\mathrm{Re}\,s=1\). Landau 1911
(\emph{Handbuch der Lehre von der Verteilung der Primzahlen})
deduced the bound on \(M_{1}(N)\). The quantitative form follows from
the de la Vall\'ee-Poussin zero-free region
\(\sigma\geq 1-c/\log(|t|+3)\).
\end{proof}

\begin{proposition}[Mertens bound, conditional on RH]
\label{v4-arith:w7-b:prop:mertens-cond}
\ClaimStatusProvedElsewhere. Under RH:
\begin{align}
\label{v4-arith:w7-b:eq:M-RH}
M(N)&\;=\;O(\sqrt{N}\log^{2}N),\\
\label{v4-arith:w7-b:eq:M1-RH}
M_{1}(N)&\;=\;O(N^{-1/2}\log^{2}N).
\end{align}
\end{proposition}

\begin{proof}
Titchmarsh \S14.27; the classical conditional bound via the explicit
formula for \(M(N)\) in terms of zeros of \(\zeta\). The bound on
\(M_{1}(N)\) follows by partial summation from the bound on \(M(N)\).
\end{proof}

\begin{remark}[Mertens bounds do not give the weighted HFD estimate]
\label{v4-arith:w7-b:rem:no-feedback}
The scalar estimates \(M_{1}(N)=o(1)\) and
\(M_{1}(N)=O(N^{-1/2}\log^{2}N)\) do not control the norm in
\eqref{v4-arith:w7-b:eq:dN-Weil-link}. That norm asks for weighted
approximation to \(1/\zeta(1/2+it)\) on the critical line. The
Baez-Duarte--Balazard--Landreau--Saias \(1/\log N\) lower obstruction
is a zero-expansion statement; without a matching global critical-line
upper theorem it cannot be transferred to
\eqref{v4-arith:w7-b:eq:HFD-NB-bound}.
\end{remark}

\section{Comparison of the Complementary Class with Baez-Duarte}
\label{v4-arith:w7-b:complementary-class}

The complement
\(\mathrm{PW}^{\mathrm{hol}}_{F}\setminus\mathrm{PW}^{\mathrm{hol},\mathrm{HFD}}_{F}\)
is the locus where the HFD archimedean positivity no longer dominates
the negative core. For test functions outside HFD, the mass on
\(|t|\leq t_{\ast}\) may be measured by a compact-core
Nyman--Beurling distance. This compact-core distance is a shadow of the
global Baez-Duarte criterion, not an equivalent replacement for it.

\begin{theorem}[complementary-class comparison with Baez-Duarte]
\label{v4-arith:w7-b:thm:complement-to-BD}
\ClaimStatusProvedHere. For
\(f\in
\mathrm{PW}^{\mathrm{hol}}_{F}\setminus\mathrm{PW}^{\mathrm{hol},\mathrm{HFD}}_{F}\),
set
\begin{equation}
\label{v4-arith:w7-b:eq:complement-BD}
d_{N}^{\mathrm{low}}(f)^{2}\;:=\;\inf_{c_{k}}\int_{|t|\leq t_{\ast}}\!|\widehat{f}(t)|^{2}\cdot\!\left|\frac{1-\zeta(s)\sum_{k=1}^{N}c_k(k^{-1}-k^{-s})}{s}\right|_{s=1/2+it}^{2}dt.
\end{equation}
On \(|t|\leq t_{\ast}\),
\begin{equation}
\label{v4-arith:w7-b:eq:core-weight-comparison}
0\leq -g(t)\leq
(1/4+t_{\ast}^{2})|g(0)|\,|s|^{-2}_{s=1/2+it}.
\end{equation}
Thus the compact-core Nyman--Beurling norm is taken over a measure
which dominates the archimedean core measure. This is a comparison of
weighted error norms, not an inequality for \(A[f]\) itself; the
finite-prime-polar functional \(R[f]\) remains in
\Cref{v4-arith:w6-a:eq:bounded-reduction-star}. The compact-core norm
does not imply the global Baez-Duarte vanishing \(d_N\to0\).
\end{theorem}

\begin{proof}
By \Cref{v4-arith:w6-a:lem:decomposition},
\(W_{\infty}(f*\widetilde{f})=(B[f]-A[f])/(2\pi)\). On
\(\mathrm{HFD}^{c}\), \(B[f]<A[f]\) is possible, so the
finite-prime-polar functional remains part of the inequality. The
negative term is
\[
A[f]=\int_{|t|\leq t_{\ast}}|\widehat{f}(t)|^{2}\cdot(-g(t))dt.
\]
By \Cref{v4-arith:w6-a:lem:g-monotone} and
\Cref{v4-arith:w6-a:prop:g-boundary},
\(-g(t)=|g(0)|-\int_{0}^{|t|}g'(u)du\geq 0\) on
\(|t|\leq t_{\ast}\). Since
\(|s|^{-2}=(1/4+t^{2})^{-1}\) on \(s=1/2+it\), the bound
\eqref{v4-arith:w7-b:eq:core-weight-comparison} follows from
\(-g(t)\leq |g(0)|\) and
\(1/4+t^{2}\leq 1/4+t_{\ast}^{2}\). Substitution into
\Cref{v4-arith:w6-a:eq:bounded-reduction-star} shows where this
measure comparison enters the bounded-interval residual. The factor
\(\left|1-\zeta(s)\sum c_k(k^{-1}-k^{-s})\right|^2\) is an additional
approximation error, so no estimate for \(A[f]\) follows without a
separate approximation theorem. The implication \(d_N\to0\) is exactly
the global Nyman--Beurling--Baez-Duarte criterion and cannot be
inferred from the compact interval alone.
\end{proof}

\begin{corollary}[compact Baez-Duarte data do not remove the irreducible A-TP residual]
\label{v4-arith:w7-b:cor:irreducible-restricted-BD}
\ClaimStatusProvedHere. The residual A-TP on
\(\mathrm{PW}^{\mathrm{hol}}_{F}\setminus\mathrm{PW}^{\mathrm{hol},\mathrm{HFD}}_{F}\)
is not closed by the compact-core norm
\eqref{v4-arith:w7-b:eq:complement-BD}. Full closure would require the
global Baez-Duarte vanishing \(d_N\to0\), equivalently RH, together
with the retained \(R[f]\)-inequality of
\Cref{v4-arith:w6-a:eq:bounded-reduction-star}.
\end{corollary}

\begin{proof}
Combine \Cref{v4-arith:w7-b:thm:complement-to-BD} with the global
criterion \Cref{v4-arith:w7-b:thm:baez-duarte}.
\end{proof}

\section{Seven Obstructions to the HFD Nyman--Beurling Bound}
\label{v4-arith:w7-b:obstruction-analysis}

The estimate in
\Cref{v4-arith:w7-b:thm:HFD-NB-bound} is the open bound
\(d_{N}^{\mathrm{HFD}}(f)^{2}\leq C_{\mathrm{HFD}}(f)/\log N+
O((\log N)^{-2})\) on the HFD sub-class. The following seven
obstructions remain after the Mellin formula is put in its pole-free
form.

\begin{remark}[obstruction 1: a scalar Mertens bound is too small]
\label{v4-arith:w7-b:rem:obstruction-1}
The limit \(M_{1}(N)=\sum_{k\leq N}\mu(k)/k\to0\) is unconditional by
PNT, and the de la Vall\'ee-Poussin zero-free region gives the usual
exponential decay in \(\sqrt{\log N}\). This scalar statement is
nevertheless too small for
\eqref{v4-arith:w7-b:eq:HFD-NB-bound}: the latter is an \(L^{2}\)
approximation estimate for \(1/\zeta(1/2+it)\) against a moving
weight.
\end{remark}

\begin{remark}[obstruction 2: the \(1/\log N\) scale is only a lower obstruction]
\label{v4-arith:w7-b:rem:obstruction-2}
The \(1/\log N\) scale is forced from below by the
Baez-Duarte--Balazard--Landreau--Saias zero expansion. A proof of an
upper bound of the same scale for
\(d_N^{\mathrm{HFD}}(f)\) would have to replace the global
zero-expansion by a weighted critical-line theorem. Truncating the
spectral weight to \(|t|>t_{\ast}\) does not by itself supply such a
theorem.
\end{remark}

\begin{remark}[obstruction 3: the HFD class is empty or trivial]
\label{v4-arith:w7-b:rem:obstruction-3}
\Cref{v4-arith:w6-a:rem:HFD-nonempty} gives explicit high-frequency
Paley--Wiener functions satisfying
\eqref{v4-arith:w6-a:eq:HFD-condition}. Non-emptiness only proves that
the archimedean inequality \(W_{\infty}\ge0\) has a genuine domain; it
does not prove the weighted Nyman--Beurling estimate on that domain.
\end{remark}

\begin{remark}[obstruction 4: \(\varepsilon_{\mathrm{pol}}(N)\) conceals the real obstruction]
\label{v4-arith:w7-b:rem:obstruction-4}
After the corrected Mellin formula there is no \((s-1)^{-1}\) polar
term in \eqref{v4-arith:w7-b:eq:dN-mellin-compact}. The obstruction is
the weighted critical-line norm in
\eqref{v4-arith:w7-b:eq:dN-Weil-link}, not a finite-dimensional polar
remainder.
\end{remark}

\begin{remark}[obstruction 5: the cross-reduction identity \eqref{v4-arith:w7-b:eq:dN-Weil-link} is circular]
\label{v4-arith:w7-b:rem:obstruction-5}
The identity \eqref{v4-arith:w7-b:eq:dN-Weil-link} is a
Mellin--Plancherel identity and contains no zero-square residue
formula. It is independent of the Weil distribution, but independence
of identities is weaker than a bound. The missing step is precisely a
uniform estimate for the displayed integral.
\end{remark}

\begin{remark}[obstruction 6: disjoint sources do not prove the missing estimate]
\label{v4-arith:w7-b:rem:obstruction-6}
The Nyman--Beurling construction uses Nyman 1950, Beurling 1955, and
Baez--Duarte 2003, a line distinct from Rankin--Selberg. Distinct
lineage does not imply
\eqref{v4-arith:w7-b:eq:HFD-NB-bound}; the required weighted estimate
remains an analytic input.
\end{remark}

\begin{remark}[obstruction 7: closing A-TP on \(\mathrm{HFD}^{c}\) via Nyman--Beurling]
\label{v4-arith:w7-b:rem:obstruction-7}
The compact-core norm \eqref{v4-arith:w7-b:eq:complement-BD} measures
only the archimedean negative interval. Closing A-TP on
\(\mathrm{HFD}^{c}\) still requires the \(R[f]\)-term in
\Cref{v4-arith:w6-a:eq:bounded-reduction-star}; replacing it by the
global Nyman--Beurling vanishing is exactly the RH-equivalent
Baez-Duarte criterion.
\end{remark}

\section{The Irreducible Obstruction, Named in the Nyman--Beurling Framework}
\label{v4-arith:w7-b:irreducible-NB}

\begin{theorem}[irreducible A-TP residual, Nyman--Beurling form]
\label{v4-arith:w7-b:thm:irreducible-NB}
\ClaimStatusProvedHere. The Nyman--Beurling formulation identifies the
same residual in two forms: the global Baez-Duarte vanishing
\(d_N\to0\), which is RH-equivalent, and the compact-core norm
\eqref{v4-arith:w7-b:eq:complement-BD}, which measures only the
negative archimedean interval. The latter does not close A-TP without
the \(R[f]\)-inequality, and the HFD weighted estimate
\Cref{v4-arith:w7-b:thm:HFD-NB-bound} remains an additional analytic
input.
\end{theorem}

\begin{proof}
Assemble \Cref{v4-arith:w7-b:thm:complement-to-BD},
\Cref{v4-arith:w7-b:cor:irreducible-restricted-BD}, and the corrected
form of \Cref{v4-arith:w7-b:thm:HFD-NB-bound}.
\end{proof}

\begin{remark}[comparison with the bounded-interval residual inequality]
\label{v4-arith:w7-b:rem:compare-w6}
The bounded-interval reduction (\Cref{v4-arith:w6-a:ATP-direct-obstruction}) reduced A-TP to a
bounded-interval domination \(B[f]\leq A[f]+2\pi R[f]\), with the
high-frequency archimedean tail controlled by the compact core and the
residual. The Nyman--Beurling
form gives a compact-core distance for the negative interval and a
global distance equivalent to RH. Neither distance removes the
\(R[f]\)-term from the bounded-interval inequality. The
Nyman--Beurling construction supplies exactly the following data:
\begin{enumerate}
\item the exact Mellin normal form for the weighted critical-line
estimate and an \(L^{2}(0,1)\) distance;
\item the precise weighted estimate whose proof would be required for
HFD closure;
\item a reformulation of the irreducible residual as a concrete
Baez-Duarte distance problem.
\end{enumerate}
It adds no closure of A-TP beyond the \(R[f]\)-coupled sub-classes. The
elementary ceiling is the same ceiling seen in Mellin coordinates.
\end{remark}

\section{Mellin Boundary}
\label{v4-arith:w7-b:summary}

\begin{enumerate}
\item The Nyman--Beurling generators
\(\rho_{\alpha}(x)=\{\alpha/x\}-\alpha\{1/x\}\) have Mellin transform
\((\alpha-\alpha^s)\zeta(s)/s\)
(\Cref{v4-arith:w7-b:lem:mellin-rho}); A-TP is equivalent to the
Nyman--Beurling closure \(\chi\in H^{2}(0,1)\) and to the Baez-Duarte
vanishing \(d_{N}\to 0\)
(\Cref{v4-arith:w7-b:thm:translation-ATP-NB}).
\item The HFD-restricted Nyman--Beurling distance
\(d_{N}^{\mathrm{HFD}}(f)\)
(\Cref{v4-arith:w7-b:def:dN-HFD}) is governed by the exact
Mellin-Plancherel integral \eqref{v4-arith:w7-b:eq:dN-Weil-link}. The
proposed \(O(1/\log N)\) HFD bound is a separate weighted estimate
(\Cref{v4-arith:w7-b:thm:HFD-NB-bound}).
\item On \(\mathrm{HFD}^{c}\), the compact-core Baez-Duarte norm
controls the negative archimedean interval but does not replace the
global Baez-Duarte criterion or the retained \(R[f]\)-inequality
(\Cref{v4-arith:w7-b:thm:complement-to-BD},
\Cref{v4-arith:w7-b:cor:irreducible-restricted-BD}).
\item The obstruction-theoretic comparison changes the form of the HFD
Nyman--Beurling construction: PNT-strength Mertens estimates are background,
not a proof of the weighted critical-line estimate.
\item The Nyman--Beurling elementary ceiling coincides with the
bounded-interval elementary ceiling. It gives a Mellin reformulation
of the residual, not an HFD theorem for A-TP
(\Cref{v4-arith:w7-b:thm:irreducible-NB}).
\end{enumerate}

\begin{remark}[the Beilinson principle on the Nyman--Beurling comparison]
\label{v4-arith:w7-b:rem:beilinson-closing}
The Nyman--Beurling comparison proves neither RH nor an independent
HFD closure. Its mathematical content is the irreducible residual
recast in \(L^{2}(0,1)\) language and the exact weighted estimate which
would imply the missing HFD bound.
\end{remark}

\begin{remark}[cross-volume dependency]
\label{v4-arith:w7-b:rem:cross-volume}
The inputs are Chapter~\ref{v4-arith:w5-a:chapter}
(Rankin--Selberg kernel obstruction, A-TP
conjecture~\ref{v4-arith:w5-a:conj:archimedean-tate-positivity}),
Chapter~\ref{v4-arith:w6-a:ATP-direct-obstruction} (HFD sub-class,
bounded-interval reduction), and PNT
(Hadamard 1896, de la Vall\'ee-Poussin 1896, Landau 1911). They are
independent of the de Branges / Hermite--Biehler analysis of
Chapter~\ref{v4-arith:w5-j:VBstar-spectral-positivity}, the
\(\mathrm{R}_{\mathrm{analytic}}\) density-of-zeros analysis of
Chapter~\ref{v4-arith:w5-i:chapter}, or the Positselski four
sub-items of
Chapter~\ref{v4-arith:w5-d:chapter}. The Nyman--Beurling construction is
orthogonal to these.
\end{remark}
